\section{Algebra and split computations}\label{sec:algebra}

This section separates the algebraic work from the later geometry. We first
represent a word by its finite-monoid value. We then show how a sound split
extends graph tests from a skipped class of episodes to a larger class,
and how affine corrections can be removed from sequence computations.
The final lemma supplies the correction that lets a split omit one interval
product.

We may assume that $\Sigma$ is nonempty: otherwise its only languages
are $\varnothing$ and $\{\eps\}$, both of generalized star height zero.
Fix a morphism
\[
 T:\Sigma^*\longrightarrow M
\]
into a finite monoid, and put $m=|M|$ and $N=m!$.
It suffices to construct bounded-height expressions for the fibers of
$T$. Indeed, a deterministic finite automaton supplies such a morphism
by assigning to a word its transformation of the state set, and its
accepted language is a finite union of fibers.

Letters occupy unit intervals with integer endpoints. For an available
interval $[a,b)$, with $a\le b$, write $T_{a,b}$ for the product of its
letters; in particular $T_{a,a}=1_M$. Products and actions are written
in reading order, so $T_{a,c}=T_{a,b}T_{b,c}$ when $a\le b\le c$.
Let $V=\F_2^M$ have basis $(e_c)_{c\in M}$. Right multiplication in $M$ induces the
linear right action $e_c d=e_{cd}$, extended to all of $V$.
The image $e_{1_M}d=e_d$ distinguishes every $d\in M$.
We will also use other finite-dimensional vector spaces over $\F_2$,
with linear right actions written $v\rho(d)$.

One elementary consequence of finiteness will be used when boundaries
move through periodic words:
\begin{equation}\label{eq:monoid-periodicity}
 c^{h+N}=c^h\qquad(c\in M,\ h\ge m).
\end{equation}
To see this, choose $0\le a<b\le m$ with $c^a=c^b$.
Multiplication by further powers gives period $b-a$ from exponent $a$
onward, and $b-a$ divides $N$. This also covers $m=N=1$.

\subsection{Splitting a sequence of episodes}

A \emph{prefix code} is a set $E\subseteq\Sigma^+$ in which no word is
a proper prefix of another. Every word in $E^*$ has a unique
factorization into elements of $E$: compare the first factors, which
must coincide by the prefix condition, and continue. We call the
elements of $E$ \emph{episodes}.

Suppose each episode $e$ carries a deterministic update $G_e$ of a
fixed finite state set $S$. A \emph{graph language} from $x$ to $y$ on
$D^*$, where $D\subseteq E$, consists of those words in $D^*$ whose
successive episode updates send $x$ to $y$. The empty word carries the
identity update. Throughout, a graph language includes its indicated
domain condition; it accepts no word outside $D^*$.

The next lemma receives two independently defined languages for the factors
of a split. It assumes universal soundness of their concatenations but asks
for a successful split only outside the skipped class. In particular,
completeness does not license either side to trust the other side's guesses.
The common update on an episode must be fixed before either guess is made.

\begin{lemma}[Split principle]\label{lem:split}
Let $D'\subseteq D\subseteq E$, where $E$ is a prefix code, and let
$(G_e)_{e\in D}$ be deterministic updates of a finite set $S$.
For $x,y\in S$, suppose languages $P_x,Q_y$ satisfy the following
conditions.
\begin{enumerate}[label=\textup{(\roman*)}]
 \item If a word in $D^*$ has a prefix $uv$ with $u\in P_x$ and
 $v\in Q_y$, then its first episode $e$ exists, $uv=e$, and
 $G_e(x)=y$.
 \item For every $e\in D\setminus D'$ and every $x\in S$,
 $e\in P_xQ_{G_e(x)}$.
\end{enumerate}
Let $W_{x,y}$ be the graph languages for the same updates on $(D')^*$.
Then the graph language from $x$ to $y$ on $D^*$ is
\begin{equation}\label{eq:split-expression}
\begin{split}
D^*\cap\Bigg[W_{x,y}\ \cup\ &
 \left(\bigcup_{a\in S}W_{x,a}P_a\right)
 \left(\bigcup_{b,c\in S}Q_bW_{b,c}P_c\right)^*\\[-2pt]
 &\hspace{22mm}\cdot
 \left(\bigcup_{f\in S}Q_fW_{f,y}\right)\Bigg].
\end{split}
\end{equation}
If $D$, all $P_x$, and all $Q_y$ have height at most one, and all
$W_{x,y}$ have height at most $b\ge0$, then these new graph languages
have height at most $\max\{2,b+1\}$.
\end{lemma}

\begin{proof}
\emph{Soundness of the expression.}
Consider a factorization accepted by the expression on the right of
\eqref{eq:split-expression}. The outer intersection puts the word in
$D^*$, so it has its true episode boundaries. A $W$-factor starting
at such a boundary consumes a sequence of complete episodes:
its factorization into members of $D'\subseteq E$ agrees with the
ambient factorization by the prefix-code property.

In the second term, pair each $P$-factor with the following $Q$-factor.
Condition~\textup{(i)} says that this pair consumes exactly the next
episode and has the indicated input and output states. Induction
through the alternating $W$-factors and these pairs proves both the
boundary alignment and the required composite update. This includes
the case in which the central star contributes no factors, leaving
one explicit split. The first term $W_{x,y}$ is already correct.

\emph{Completeness and height.}
Conversely, take a word in $D^*$ with the required update. Split each
episode outside $D'$ using condition~\textup{(ii)}, and absorb each
intervening sequence of $D'$-episodes into its graph language $W$.
This gives the second term, unless there are no explicit splits, in
which case it gives $W_{x,y}$. Thus the displayed language is exact.

For the height estimate, $D^*$ has height at most two. The central
star has height at most $1+\max\{1,b\}$, and all other operations are
finite unions, concatenations, or intersections. The stated bound
follows.
\end{proof}

\subsection{Recovering linear parts}

The split identity has reduced the language problem to episode updates.
We are free to add an episode-dependent affine correction if we can later
recover the desired linear action. The next identity does that recovery
on the same complete word in both tests. It is therefore valid after
arbitrarily many episode updates, not merely after one.

\begin{lemma}[Affine recovery]\label{lem:affine-recovery}
Let $U$ be a finite-dimensional vector space over $\F_2$ and let
$D$ be a subset of a prefix code. Suppose the update for $e\in D$ is
\[
 G_e(v)=vL_e+\sigma_e,
\]
where $L_e$ is linear and $\sigma_e\in U$, both determined by the
episode. If $R_{v,z}$ are the graph languages of the affine composite
on $D^*$, then the graph language of the product of the linear parts,
from $v$ to $y$, is
\[
 \bigcup_{f\in U}\bigl(R_{v,y+f}\cap R_{0,f}\bigr).
\]
In particular this recovery does not increase a common height bound.
\end{lemma}

\begin{proof}
Applying $v\mapsto vL+c$ and then $v\mapsto vA+d$ gives
$v\mapsto vLA+cA+d$. Hence the composite on any word of $D^*$ has
the form $v\mapsto vL+c$, where $L$ is the product of the episode
linear parts in reading order. Its output at zero is $c$.
The displayed intersection says that the two outputs are $y+f$ and
$f$, which is equivalent to $vL=y$. The union ranges over the finite
set of possible $f$. Only finite Boolean operations are used.
\end{proof}

\subsection{An affine correction that hides one interval product}

At an interior cut, neither factor can read the product across the
cut. The following finite lemma provides a correction for which one
interval product can be omitted. The omitted interval may depend on
the input state and the full tuple, but the correction itself does
not depend on the input state.

\begin{lemma}[Affine-line constancy]\label{lem:affine-line}
Let $U$ be a finite-dimensional vector space over $\F_2$ and let
$\rho$ be a linear right action of $M$ on $U$. There are an integer
$l\ge1$ and a function $\beta:M^l\to U$ such that, for every
$v\in U$ and every $\mathbf d=(d_1,\ldots,d_l)\in M^l$, some
$j\in\{1,\ldots,l\}$ makes the function
\[
 a\longmapsto
 v\rho(d_1\cdots d_{j-1}a d_{j+1}\cdots d_l)
 +\beta(d_1,\ldots,d_{j-1},a,d_{j+1},\ldots,d_l)
\]
constant on $M$.
\end{lemma}

\begin{proof}
We first bound the probability of failure for one state and tuple, then
show that all these failures can be avoided simultaneously. The length
$l$ is chosen only after both estimates are available.

Write $q_U=|U|$. If $m=1$ or $q_U=1$, take $l=1$ and $\beta=0$.
We therefore assume $m,q_U>1$.
For an integer $l$ to be chosen, choose the values of $\beta$
independently and uniformly in $U$ at all points of $M^l$.
For each $(v,\mathbf d)$, let $\mathcal B_{v,\mathbf d}$ be the event
that no coordinate line through $\mathbf d$ has the asserted
constancy.

Condition on $\beta(\mathbf d)$. Constancy along a specified line
then prescribes the values of $\beta$ at its other $m-1$ points, so
has probability $q_U^{1-m}$. Different coordinate lines through
$\mathbf d$ have disjoint sets of noncentral points. Their constancy
events are therefore independent under this conditioning, and the
conditional probability of failure does not depend on the center
value. Consequently
\[
 \Pr(\mathcal B_{v,\mathbf d})
 =p_l:=(1-q_U^{1-m})^l.
\]

An event $\mathcal B_{v,\mathbf d}$ uses only the random values at
points of Hamming distance at most one from $\mathbf d$. It is thus
independent of the joint family of events whose centers have distance
greater than two from $\mathbf d$. The number of potentially
dependent events is at most
\[
 q_U\left(1+l(m-1)+\binom l2(m-1)^2\right)-1
 \ \le\ D_l:=q_U(1+lm+l^2m^2).
\]
This count includes all possible input vectors at each nearby center.

We use a local-lemma argument \citep[Section~2, pp.~616--617]{ErdosLovasz1975},
recalling the required finite probability estimate with its proof.
Suppose a finite family of events has a dependency graph of maximum
degree at most $D\ge1$: each event is independent of the joint family
of its nonneighbors. If every event has probability at most
$u(1-u)^D$, where $0<u<1$, then their simultaneous nonoccurrence
has positive probability.

For completeness, induct on the size of a conditioning set to show
that its nonoccurrence has positive probability and that the
conditional probability of any other event is at most $u$.
The assertion with an empty conditioning set follows from the
probability bound. Positivity for a set of size $k$ follows by
removing one member and applying the conditional bound for size
$k-1$. To prove the conditional bound for an event $A$ and a set
$S$ of size $k$ not containing it, partition $S$ into its neighbors
$S_1$ and nonneighbors $S_0$. If $S_1$ is empty, joint independence
gives the unconditional bound. Otherwise,
\[
 \Pr\left(A\,\middle|\,\bigcap_{B\in S}B^c\right)
 \le
 \frac{\Pr\left(A\,\middle|\,\bigcap_{B\in S_0}B^c\right)}
 {\Pr\left(\bigcap_{B\in S_1}B^c\,\middle|\,
                    \bigcap_{B\in S_0}B^c\right)}
 \le \frac{\Pr(A)}{(1-u)^{|S_1|}}
 \le u.
\]
The denominator is bounded by revealing the members of $S_1$
successively: every conditional set then has size at most $k-1$,
so the induction applies. The numerator uses joint independence
from $S_0$. This completes the induction and proves the estimate.

Apply the estimate with $D=D_l$ and $u=1/(2D_l)$. Bernoulli's
inequality gives
\[
 u(1-u)^{D_l}
 =\frac1{2D_l}\left(1-\frac1{2D_l}\right)^{D_l}
 \ge\frac1{4D_l}.
\]
The quantity $p_l$ decays exponentially in $l$, whereas $D_l$ grows
quadratically. Choosing $l$ large enough makes $p_l\le1/(4D_l)$.
There is therefore a choice of $\beta$ for which no failure event
occurs. Fix such a choice.
\end{proof}

\begin{remark}[Uniformity and finite choice]\label{rem:finite-affine-choice}
The choices of $l$ and $\beta$ precede the state $v$ and tuple $\mathbf d$.
Only the successful coordinate $j$ may depend on them. The proof is also
an effective finite existence argument: increase $l$ until its displayed
probability inequality holds, then enumerate the finitely many functions
$M^l\to U$ and test the stated property for every state and tuple. At least
one function passes. This procedure asserts no useful running-time bound.
\end{remark}

To apply the lemma to a word segment, choose ordered checkpoints
$B_0<\cdots<B_l$ and set $d_j=T_{B_{j-1},B_j}$.
Process the segment by $\rho$, then add $\beta(\mathbf d)$ at $B_l$.
If a cut lies in interval $j$, the prefix can transmit the vector
$v$ at $B_0$ and the preceding products $d_1,\ldots,d_{j-1}$.
The suffix supplies $d_{j+1},\ldots,d_l$ and checks constancy on the
missing-coordinate line. When that check holds, it obtains the
correct vector at $B_l$ without knowing $d_j$. For every actual
tuple and input vector at least one interval permits this procedure.
Because the added vector is a function of the actual tuple alone,
the update of the whole segment is affine and remains eligible for
Lemma~\ref{lem:affine-recovery}.
